% ---------------------------------------------------------------------
%  Chapitre 6 — Conclusion (version aperçu, rédigée)
% ---------------------------------------------------------------------
\chapter{Conclusion}
\label{chap:conclusion}

L'objectif fixé en introduction n'était pas de découvrir une stratégie
gagnante, mais de construire un protocole capable de juger n'importe
laquelle. Ce protocole est construit, éprouvé, et il a rendu son
verdict~: aucune des dix stratégies d'analyse technique ne démontre
d'avantage réel une fois ses conditions d'application vérifiées.

Le résultat central mérite d'être rappelé. Le test standard déclarait
deux stratégies gagnantes. La vérification de la seule condition
d'indépendance a renversé les deux verdicts, l'un se révélant un faux
positif né de la répétition d'une même information sur des milliers de
trades, l'autre un avantage réel mais trop épisodique pour être retenu.
Dans le même temps, le témoin construit pour être gagnant a été validé
sans hésitation à chaque étage. Un outil de validation se juge à sa
capacité d'attraper ses propres erreurs tout en reconnaissant le vrai~:
la démonstration est faite.

Une idée traverse l'ensemble de ce travail. La dépendance des données y
a joué les deux rôles. Simultanée, portée par le facteur marché qui
explique près des trois quarts des mouvements, elle était la nuisance
qui faussait les tests, et il a fallu la mesurer puis la neutraliser
pour juger proprement. Retardée, elle aurait été la seule ressource
qu'une stratégie puisse exploiter, et le marché n'en contient presque
pas, conformément à l'efficience faible. Neutraliser l'une, chercher
l'autre~: c'est la même statistique, vue des deux côtés du même objet.

Le protocole est désormais prêt à juger les prochains signaux de
l'agent, en commençant par leur validation temporelle.
